\documentclass[11pt]{article}

\usepackage[a4paper,margin=1in]{geometry}
\usepackage{amsmath,amssymb,amsthm,mathtools}
\usepackage{graphicx}
\usepackage{float}
\usepackage{microtype}
\usepackage{xcolor}
\usepackage{enumitem}
\usepackage[colorlinks=true,linkcolor=blue!55!black,citecolor=blue!55!black,urlcolor=blue!60!black]{hyperref}

\newtheorem{theorem}{Theorem}
\newtheorem{lemma}{Lemma}
\newcommand{\one}{\mathbf 1}
\newcommand{\C}{\mathbb C}
\newcommand{\R}{\mathcal R}
\newcommand{\spb}{\operatorname{s}}
\newcommand{\ran}{\operatorname{ran}}

\title{Eventual Norm Continuity Does Not Preserve the\\Spectral-Gap Necessity for Eventual Domination}
\author{Automated mathematics run \texttt{fa\_banach\_001}\\Lane 0 --- candidate for human review}
\date{11 August 2026}

\begin{document}
\maketitle

\begin{center}
\fbox{\parbox{0.91\textwidth}{\textbf{Status.}
Candidate counterexample, likely valid.  This gives a full negative answer to
the first explicit question in source Remark~3.3: implication
\textup{(i)}$\Rightarrow$\textup{(iii)} of Theorem~3.1 fails when analyticity
is replaced by eventual norm continuity.  The separate question in the same
remark about removing one smoothing assumption is not resolved here.}}
\end{center}

\section{Source question and answer}

Gl\"uck and Mugnolo characterize individual eventual domination between two
distinct real analytic semigroups under spectral, smoothing, and eventual
positivity hypotheses \cite[Theorem~3.1]{GlueckMugnolo2021}.  Their assertion
\textup{(i)} says that for every $0<f$ the orbit of $B$ eventually dominates
the nonnegative orbit of $A$; assertion \textup{(iii)} is the strict spectral
gap $\spb(B)>\spb(A)$.  Immediately after the proof they ask whether
\textup{(i)}$\Rightarrow$\textup{(iii)} remains true when analyticity is
weakened, for instance to eventual norm continuity.

\begin{figure}[H]
  \centering
  \includegraphics[width=\textwidth]{figures/open_question_crop.png}
  \caption{Source PDF page~9, Remark~3.3.  The first displayed open question is
  answered negatively below.}
\end{figure}

The answer is no.  The mechanism is finite-time erasure: two different
semigroups may have exactly the same tail without being the same semigroup.
We retain every other hypothesis of the source theorem, including nonempty
spectrum, pole conditions, both smoothing inclusions, and the stated
positivity assumptions.

\section{A nilpotent shift with empty generator spectrum}

Let $H=L^2(0,1;\C)$ and define the left-shift semigroup $(S(t))_{t\geq0}$ by
\[
 (S(t)g)(x)=
 \begin{cases}
   g(x+t),&x+t\leq1,\\
   0,&x+t>1.
 \end{cases}
\]
It is a real $C_0$-semigroup and $S(t)=0$ for $t\geq1$.  Its generator is
\[
 Gg=g',\qquad
 D(G)=\{g\in H^1(0,1):g(1)=0\}.
\]

\begin{lemma}\label{lem:empty-spectrum}
The spectrum of $G$ is empty.
\end{lemma}

\begin{proof}
For every $\lambda\in\C$ and $h\in H$, set
\[
 (K_\lambda h)(x)=\int_x^1 e^{\lambda(x-s)}h(s)\,ds.
\]
This bounded Volterra operator maps $H$ into $D(G)$.  Differentiation gives
$(K_\lambda h)'=\lambda K_\lambda h-h$, so
$(\lambda-G)K_\lambda h=h$.  Conversely, integration of the scalar
differential equation with terminal value $g(1)=0$ gives
$K_\lambda(\lambda-G)g=g$ for $g\in D(G)$.  Hence
$K_\lambda=\R(\lambda,G)$ for every complex $\lambda$, proving
$\sigma(G)=\varnothing$.
\end{proof}

\section{Counterexample}

On the Banach lattice $E=L^2(0,1;\C)$, put $u=\one$ and define
\[
 \varphi(f)=\int_0^1 f(s)\,ds,
 \qquad Pf=\varphi(f)\one,
 \qquad Q=I-P,
 \qquad F=\ker\varphi.
\]
Thus $P$ is a real rank-one projection, $E=\operatorname{span}\{\one\}
\oplus F$, and $F$ is a real-invariant separable infinite-dimensional Hilbert
space.  Choose a complex-linear unitary $J:H\to F$ which maps real functions
to real functions; such a $J$ is obtained by choosing real orthonormal bases
and complexifying.  Write
\[
 W(t)=J S(t)J^{-1}\quad\text{on }F.
\]
Finally, define two operator families on $E$ by
\begin{equation}\label{eq:semigroups}
 T_A(t)=P+W(t)Q,
 \qquad
 T_B(t)=P+W(2t)Q.
\end{equation}

\begin{theorem}\label{thm:counterexample}
The families in \eqref{eq:semigroups} are distinct real $C_0$-semigroups that
satisfy every assumption of source Theorem~3.1 after ``analytic'' is replaced
by ``eventually norm continuous.''  Assertion \textup{(i)} of that theorem
holds, but
\[
 \spb(A)=\spb(B)=0,
\]
so assertion \textup{(iii)} fails.
\end{theorem}

\begin{proof}
\medskip\noindent\textbf{Semigroup law, reality, and distinctness.}
Relative to $E=\operatorname{span}\{\one\}\oplus F$, the two families are
unitarily similar to
\[
 1\oplus S(t)\quad\text{and}\quad 1\oplus S(2t),
\]
respectively.  They are therefore real $C_0$-semigroups.  They are distinct:
if $g=\one\in H$ and $h=Jg\in F$, then at $t=\tfrac14$ their complementary
components are $J\one_{(0,3/4)}$ and $J\one_{(0,1/2)}$.

\medskip\noindent\textbf{Eventual norm continuity and smoothing.}
Nilpotence of $S$ gives
\[
 T_A(t)=P\quad(t\geq1),
 \qquad
 T_B(t)=P\quad(t\geq\tfrac12).
\]
Both semigroups are thus constant in operator norm for all sufficiently large
times.  Moreover
\[
 E_u=L^\infty(0,1),
 \qquad T_A(1)E=T_B(1)E=PE=\operatorname{span}\{\one\}\subseteq E_u.
\]
This verifies both smoothing inclusions with the common time $t_0=1$.

\medskip\noindent\textbf{Spectra and peripheral poles.}
The generators are, under the same decomposition,
\[
 A=0\oplus JGJ^{-1},
 \qquad B=0\oplus 2JGJ^{-1}.
\]
Lemma~\ref{lem:empty-spectrum} and its scaled version yield
$\sigma(A)=\sigma(B)=\{0\}$.  More precisely, near zero,
\begin{align*}
 \R(\lambda,A)
   &=\lambda^{-1}P+J\R(\lambda,G)J^{-1}Q,\\
 \R(\lambda,B)
   &=\lambda^{-1}P+J\R(\lambda,2G)J^{-1}Q.
\end{align*}
The second terms are operator-valued entire functions of $\lambda$.  Thus
$0$ is a simple pole for both generators, their peripheral spectra are
$\{0\}$, and both spectral bounds equal zero.  In particular, the
construction does not evade the source's nonempty-spectrum condition.

\medskip\noindent\textbf{Positivity and domination.}
If $0<f\in E$, then strict positivity of integration gives
\[
 Pf=\varphi(f)u\gg_u0.
\]
Hence $(T_A(t))$ is uniformly eventually positive and $(T_B(t))$ is uniformly
eventually strongly positive with respect to $u$; these are stronger than the
two source positivity assumptions.  Finally, for every $0<f\in E$ and every
$t\geq1$,
\[
 T_B(t)f=T_A(t)f=Pf\geq0.
\]
Source assertion \textup{(i)} therefore holds, even with a time independent
of $f$, while the strict inequality in \textup{(iii)} fails.  This completes
the hypothesis crosswalk and the counterexample.
\end{proof}

\section{Why analyticity is the exact fault line}

In the equal-spectral-bound case, the published proof eventually derives
$T_B(t)f=T_A(t)f$ for all sufficiently large $t$ and every $0<f$.  Analytic
unique continuation then extends this equality to every positive time,
contradicting distinctness.  Eventual norm continuity has no analogous
backward-uniqueness property.  In the example above the complementary
subspace $F$ is erased in finite time, after which two different histories
are indistinguishable.

The semigroups above cannot be analytic.  Indeed, $T_A(t)=P$ for $t\geq1$ has
the nontrivial kernel $F$, whereas every positive-time operator of an analytic
$C_0$-semigroup is injective: if $T(t_0)x=0$, the semigroup law makes the
analytic orbit vanish on a tail, the identity theorem makes it vanish on
$(0,\infty)$, and strong continuity at zero gives $x=0$.

\section{Deep-upgrade attempt on the second source question}

Remark~3.3 separately asks whether the inclusion $e^{t_0B}E\subseteq E_u$ is
needed for \textup{(iii)}$\Rightarrow$\textup{(ii)}.  A later refinement by
Arora and Gl\"uck proves domination without that $B$-smoothing inclusion if
one directly assumes that for every $0<f$ there are $c_f>0$ and $t_f$ such
that
\[
 e^{tB}f\geq c_f u\qquad(t\geq t_f)
\]
\cite[Proposition~4.2]{AroraGlueck2023}.  This is stronger than individual
eventual strong positivity: in the latter definition the witness in
$e^{tB}f\gg_u0$ may depend on $t$.

The additional attack tested whether spectral convergence or a Baire
argument forces such a time-uniform lower constant.  Strong convergence in
$E$ does not control the gauge norm of $E_u$, whose positive order interior
is generally not open in the ambient norm; this is precisely where the source
proof uses $B$-smoothing.  Conversely, a moving-spike construction must obey
both the semigroup law and eventual strong positivity for every positive
initial vector, and no rigorous example survived those constraints.  Thus the
second question remains outside the scope of the present full counterexample
to the first question.

\section{Novelty and verification}

The proof was audited in nine separate checks: the nilpotent semigroup law,
the all-parameter resolvent formula, the real unitary decomposition, the two
semigroup laws, distinctness, spectrum and pole order, smoothing, positivity,
and the domination quantifiers.  No numerical approximation is used.

A bounded search on 11 August 2026 covered the run indexes, the local parsed
full-source corpus, the exact source wording, the exact arXiv id and title,
and close arXiv searches involving eventual norm continuity, eventual
equality, nilpotent shift semigroups, empty spectrum, and rank-one tails.  The
later refinement \cite{AroraGlueck2023} was also inspected directly.  No prior
answer to the first source question or matching construction was found.  The
construction is elementary, so novelty confidence is medium; mathematical
validity confidence is high pending human review.

\paragraph{Human review focus.}
Check the real-unitary identification $J$, the resolvent formula in
Lemma~\ref{lem:empty-spectrum}, the direct-sum pole calculation, and the exact
match with every hypothesis of source Theorem~3.1 under the proposed
regularity weakening.

\begin{thebibliography}{9}
\bibitem{GlueckMugnolo2021}
J.~Gl\"uck and D.~Mugnolo,
\emph{Eventual domination for linear evolution equations},
Math. Z. 297 (2021), 301--330;
\href{https://arxiv.org/abs/1912.13057}{arXiv:1912.13057v2}.

\bibitem{AroraGlueck2023}
S.~Arora and J.~Gl\"uck,
\emph{Criteria for eventual domination of operator semigroups and resolvents},
Operator Theory: Advances and Applications 292 (2023), 1--26;
\href{https://arxiv.org/abs/2204.00146}{arXiv:2204.00146v3}.
\end{thebibliography}

\end{document}
